\subsection{Weak topologies}

DEFINITION 2. --- Let F and G be two vector spaces put in duality by the bilinear form B. The coarsest topology on F that makes all the linear forms \(\mathbf{B}(., y): x \mapsto \mathbf{B}(x, y)\) continuous, where y varies in G, is called the weak topology on F defined by the duality between F and G, and we denote it by \(\sigma(\mathbf{F}, \mathbf{G})\) .

Similarly we define the weak topology \(\sigma(\mathbf{G},\mathbf{F})\) on \(\mathbf{G}\) , interchanging \(\mathbf{F}\) and \(\mathbf{G}\) in definition 1; this possibility of interchanging \(\mathbf{F}\) and \(\mathbf{G}\) applies to all the results and definitions that follow in this paragraph.

We use the adjective « weak » and the adverb « weakly » to denote properties relative to a weak topology \(\sigma(\mathrm{F},\mathrm{G})\) provided there is no possibility of confusion. We shall speak, for example, of « weak convergence » and « weakly continuous functions » etc.

When \(G \subset F^{*}\) , the notation \(\sigma(F, G)\) will always denote the weak topology defined by the duality corresponding to the restriction to \(F \times G\) of the canonical bilinear form \((x, x^{*}) \mapsto \langle x, x^{*} \rangle\) .

Without extra hypotheses on F and G, we often write \(\langle x, y \rangle\) for the value \(\mathbf{B}(x, y)\) of the bilinear form B at \((x, y)\) , provided there is no ambiguity; we shall adopt this convention in the rest of this paragraph.

A vector space F carrying a weak topology of \(\sigma(F, G)\) will be called a weak space. A weak topology \(\sigma(F, G)\) is locally convex (II, p.~26, prop. 4); more precisely, it is the inverse image of the product topology of \(\mathbf{R}^G\) by the linear mapping \(\phi: x \mapsto (\langle x, y \rangle)_{y \in G}\) of \(F\) in \(\mathbf{R}^G\) . It is defined by the set of semi-norms \(x \mapsto |\langle x, y \rangle|\) when \(y\) varies in \(G\) (II, p.~5). For every \(\alpha > 0\) , and every finite family \((y_i)_{1 \leq i \leq n}\) of points of \(G\) , let \(W(y_1, ..., y_n; \alpha)\) be the set of the \(x \in F\) such that \(|\langle x, y_i \rangle| \leqslant \alpha\) for \(1 \leq i \leq n\) ; these sets (for \(\alpha, n\) and \(y_i\) arbitrary) form a fundamental system of neighbourhoods of 0 for \(\sigma(F, G)\) . Note that \(W(y_1, ..., y_n; \alpha)\) contains that vector subspace of \(F\) , of finite codimension, which is defined by the equations \(\langle x, y_i \rangle = 0\) for \(1 \leq i \leq n\) .

PROPOSITION 2. --- The weak topology \(\sigma(F, G)\) is Hausdorff if and only if the duality between \(F\) and \(G\) is separating in \(F\) .

This is a particular case of II, p.~3, prop. 2.

PROPOSITION 3. --- Let F and G be two real vector spaces in duality. Every linear form on F, that is continuous for \(\sigma(F, G)\) , can be written as \(x \mapsto \langle x, y \rangle\) for some \(y \in G\) . The element \(y \in G\) is unique when the duality is separating in G.

For, to say that the linear form \(f\) on \(\mathbf{F}\) is continuous for \(\sigma(\mathbf{F}, \mathbf{G})\) means that there exists a finite set of points \(y_i \in \mathbf{G}\) ( \(1 \leqslant i \leqslant n\) ) such that, for all \(x\) in \(\mathbf{F}\) , \(|f(x)| \leqslant \sup_{1 \leqslant i \leqslant n} |\langle x, y_i \rangle|\) (II, p.~6, prop. 5). The \(n\) relations \(\langle x, y_i \rangle = 0\) ( \(1 \leqslant i \leqslant n\) ) imply therefore \(f(x) = 0\) , and hence (A, II, § 7.5, cor. 1), there exists a linear combination \(y = \sum_{i=1}^{n} \lambda_i y_i\) such that \(f(x) = \langle x, y \rangle\) for all \(x \in \mathbf{F}\) . The uniqueness follows from \((\mathrm{D}_{\mathrm{II}})\) .

In other words, when the duality is separating in G, and F has the topology \(\sigma(\mathrm{F}, \mathrm{G})\) , then we can identify G canonically with the dual of F for this topology (II, p.~42, def. 1).

COROLLARY 1. --- A family \((a_i)\) of points of F is total for the topology \(\sigma(\mathrm{F},\mathrm{G})\) if, and only if, for every \(y \neq 0\) in G, there exists an index \(\iota\) such that \(\langle a_{\iota}, y \rangle \neq 0\) .

For using prop. 3 and I, p.~13, th. 1, the property expresses the fact that for \(\sigma(F, G)\) no closed hyperplane contains all the \(a_{i}\) ; the corollary follows therefore from cor. 3 of II, p.~38.

COROLLARY 2. --- A family \((a_{\iota})\) of points of F is topologically independent for the topology \(\sigma(\mathrm{F}, \mathrm{G})\) , if, and only if, for every index \(\iota\), there exists an element \(b_{\iota} \in G\) such that: \(\langle a_{\iota}, b_{\iota} \rangle \neq 0\) and \(\langle a_{\kappa}, b_{\iota} \rangle = 0\) for all \(\kappa \neq \iota\) .

This means, that for all \(\iota\) , there exists a closed hyperplane in \(\sigma(F, G)\) , which contains all the \(a_{\kappa}\) with index \(\kappa \neq \iota\) but does not contain \(a_{\iota}\) .

COROLLARY 3. --- Let \(G_{1}\) and \(G_{2}\) be two vector subspaces of \(F^{*}\) , in duality with F (for the restriction of the canonical bilinear form). Then \(\sigma(F, G_{2})\) is finer than \(\sigma(F, G_{1})\) if and only if \(G_{1} \subset G_{2}\) .

The condition is obviously sufficient; conversely, if \(\sigma(F, G_2)\) is finer than \(\sigma(F, G_1)\) , then every linear form that is continuous for \(\sigma(F, G_1)\) is also continuous for \(\sigma(F, G_2)\) , hence \(G_1 \subset G_2\) by prop. 3.

COROLLARY 4. --- Let G be a vector subspace of the dual F*, of the vector space F. Then F and G are in separating duality (for the canonical bilinear form) if, and only if, G is dense in F* in the topology σ(F*, F).

This follows from cor. 1.

